% Generated draft: requires mathematical and source review.
\documentclass[11pt]{article}
\usepackage[margin=25mm]{geometry}
\usepackage{fontspec}
\setmainfont{Latin Modern Roman}
\usepackage{amsmath,amssymb,mathtools,unicode-math}
\setmathfont{Latin Modern Math}
\usepackage{enumitem,needspace,xurl,hyperref}
\setlength{\parindent}{0pt}
\setlength{\parskip}{5pt}
\setlength{\emergencystretch}{3em}
\newcommand{\N}{\mathbb{N}}
\newcommand{\Z}{\mathbb{Z}}
\newcommand{\Q}{\mathbb{Q}}
\newcommand{\R}{\mathbb{R}}
\newcommand{\C}{\mathbb{C}}
\newcommand{\F}{\mathbb{F}}
\newcommand{\A}{\mathbb{A}}
\newcommand{\PP}{\mathbb{P}}
\newcommand{\E}{\mathbb{E}}
\newcommand{\eps}{\varepsilon}
\newcommand{\dd}{\,\mathrm{d}}
\begin{document}
\section*{UnsolvedMath ID: 20000398}
\textbf{Source problem code:} catalogue entry
\subsection{Definitions and mathematical statement}
Melanie Matchett Wood asks when principles labeled (A) and (B) predict an interesting average of $\#C(\F_q)$. The record does not state (A), (B), the family of curves, averaging measure, limiting parameter or meaning of ``interesting.'' Its verification expressly calls the aid an unverified reconstruction and warns that recovering workshop notes could change the labels or normalization. Thus no precise assertion about the original principles can be recovered from the wording alone.

The attached aid works instead inside a tame structured local mass model for a finite transitive permutation group $G$ on a set $\Omega$. Let $\mathcal T$ be its specified finite set of Frobenius--inertia types, $w_t$ their nonnegative masses, and $r_t$ their rational-fiber-size statistic, with $Z=\sum_t w_t>0$. It reports the model identity
\[
\E(r)=1+\frac{\sum_{t\text{ ramified}}w_t(r_t-1)}{Z},
\]
using cancellation of the unramified contribution. In the uniform unramified permutation model that cancellation is $|G|^{-1}\sum_{g\in G}|\Omega^g|=1$, by transitivity and Burnside's lemma. The aid describes masses using centralizers and discriminants of tame Frobenius--inertia pairs $(x,y)$, but does not supply their complete definitions or formula here. The identity is conditional on those masses, not a theorem that every geometric family realizes them.

For a finite cover $C\to B$, the relation $\#C(\F_q)=\sum_{b\in B(\F_q)}\#C_b(\F_q)$ explains why first-moment predictions require only local marginal distributions, not independence between different base points. The base and cover structure are additional interpretations absent from the original. A local mean fiber size must not silently be identified with the unnormalized total point count.

Within its chosen model the aid reports exact zero bias for regular abelian monodromy, and recovery of Wood's partition formula for the natural $S_n$ action: mean one for $n\le2$, strict upward bias for $n\ge3$. It also quotes $1+q^{-1}+O_n(q^{-2})$; this cannot apply to the mean-one $n\le2$ cases, so its intended degree range must be specified. The exact partition formula, tameness restrictions and geometric mass-realization hypotheses are not supplied. These model calculations do not resolve which original principles (A),(B) apply to which curve families.
\subsection{Short English statement}
Clarify the missing principles and averaging family before predicting curve point counts. The aid gives a conditional local ramification-bias formula and explains why local first moments, without independence, can determine a global first moment.
\subsection{Sources}
\begin{itemize}
\item Melanie Matchett Wood, Question 20, AIM \emph{Arithmetic statistics over finite fields and function fields}: \url{https://aimath.org/pastworkshops/arithstatffieldproblems.pdf}. Canonical location: \texttt{aim-algebraic-number-theory-notes.json}, notes[49]. The record is reconstructed/unverified, with no clean statement, original remarks or literature. No PDF or workshop slides were consulted.
\item The supplied verification warns that principles (A),(B) and their normalization are reconstructed. The aid invokes Burnside's lemma and Wood's partition formula without full bibliographic references or the local-mass definitions. It is machine generated and not peer reviewed, with low novelty confidence.
\item Local registry: \texttt{data/unsolvedmath/problems.json}, permanent catalogue entry, exact source code \texttt{catalogue entry}.
\end{itemize}
\end{document}
